\documentclass[a4paper]{arrowhead}

\usepackage[yyyymmdd]{datetime}
\usepackage{etoolbox}
\usepackage[T1]{fontenc}
\usepackage[utf8]{inputenc}
\usepackage{multirow}
\usepackage{longtable}
\usepackage{amssymb}
\usepackage{float}
\usepackage{tikz}
\usetikzlibrary{positioning,arrows.meta,fit,calc,shapes.misc,shapes.geometric}

\renewcommand{\dateseparator}{-}
\setlength{\emergencystretch}{3em}
\lstset{basicstyle=\ttfamily\footnotesize,breaklines=true,columns=fullflexible,frame=single,rulecolor=\color{gray!50},backgroundcolor=\color{gray!6}}
\tikzset{
  ahsys/.style={draw=ArrowheadBlue,very thick,fill=ArrowheadBlue!6,rounded corners=2pt,minimum width=3.2cm,minimum height=1.1cm,align=center,font=\small\bfseries},
  ahext/.style={draw=ArrowheadGrey,thick,dashed,fill=gray!5,rounded corners=2pt,minimum width=3.2cm,minimum height=1.0cm,align=center,font=\small},
  ahsrv/.style={draw=ArrowheadPurple,thick,fill=ArrowheadPurple!8,rounded corners=6pt,minimum height=0.55cm,inner xsep=4pt,align=center,font=\scriptsize\ttfamily},
  ahcsm/.style={draw=ArrowheadTeal,thick,fill=ArrowheadTeal!8,rounded corners=6pt,minimum height=0.55cm,inner xsep=4pt,align=center,font=\scriptsize\ttfamily},
  ahdb/.style={draw=ArrowheadOrange,thick,fill=ArrowheadOrange!8,cylinder,shape border rotate=90,aspect=0.25,minimum height=0.9cm,align=center,font=\scriptsize},
  ahflow/.style={-{Stealth[length=2.2mm]},thick,ArrowheadGrey},
  ahlabel/.style={font=\scriptsize,fill=white,inner sep=1pt,align=center},
}

%% Special references
\newcommand{\fref}[1]{{\textcolor{ArrowheadBlue}{\hyperref[sec:functions:#1]{#1}}}}
\newcommand{\mref}[1]{{\textcolor{ArrowheadPurple}{\hyperref[sec:model:#1]{#1}}}}
\newcommand{\sref}[1]{{\textcolor{ArrowheadPurple}{#1}}}
\newrobustcmd\msubsection[2]{
  \addtocounter{subsection}{1}
  \addcontentsline{toc}{subsection}{\protect\numberline{\thesubsection}#1 \textcolor{ArrowheadPurple}{#2}}
  \subsection*{\thesubsection\quad#1 \textcolor{ArrowheadPurple}{#2}}
  \label{sec:model:#2}
}

\begin{document}

%% Arrowhead Document Properties
\ArrowheadTitle{WorkflowExecutor}

\ArrowheadType{System Design Description}
\ArrowheadTypeShort{SysDD}
\ArrowheadVersion{5.4.0}
\ArrowheadDate{2026-11-24}
\ArrowheadAuthor{Documentation maintenance}
\ArrowheadStatus{REVIEW DRAFT}
\ArrowheadContact{Repository maintainers}
\ArrowheadFooter{\href{www.arrowhead.eu}{www.arrowhead.eu}}
\ArrowheadSetup
%%

%% Front Page
\begin{center}
  \vspace*{1cm}
  \huge{\arrowtitle}

  \vspace*{0.2cm}
  \LARGE{\arrowtype}
  \vspace*{1cm}

  
  \vspace*{\fill}

  

  \vspace*{1cm}
  \vspace*{\fill}

  % Front Page Abstract
  \begin{abstract}
    This document describes how the WorkflowExecutor is implemented in the Java reference implementation: its internal components, the realisation of its functional and non-functional properties, its configuration and the implementation status of the other languages. Documentation edition 5.4.0.
  \end{abstract}

  \vspace*{1cm}
\end{center}

\newpage
%%

%% Table of Contents
\tableofcontents
\newpage
%%

\section{Overview}
\label{sec:overview}
This document describes the design of the WorkflowExecutor, which runs the steps of workflow sessions: it has providers orchestrated for itself, invokes them through equipment adapters and reports the outcome to the WorkflowManager. The abstract system is described in the
SysD of the same name; its produced services are specified by SDs and IDDs in this documentation set.

\begin{figure}[ht!]
  \centering
  \resizebox{\textwidth}{!}{
\begin{tikzpicture}[node distance=4mm]
\node[ahsys,minimum height=6.75cm,minimum width=3.6cm] (sys) at (0,0) {WorkflowExecutor\\[-1pt]{\normalfont\scriptsize $\ll$system$\gg$}};
\node[ahsrv,anchor=west] (p0) at (3.4,3.00) {workflowExecution};
\draw[thick,ArrowheadPurple] (1.8,3.00) -- (p0.west);
\fill[ArrowheadPurple] (1.8,3.00) circle (1.4pt);
\node[ahsrv,anchor=west] (p1) at (3.4,2.25) {generalManagement};
\draw[thick,ArrowheadPurple] (1.8,2.25) -- (p1.west);
\fill[ArrowheadPurple] (1.8,2.25) circle (1.4pt);
\node[ahcsm,anchor=east] (c0) at (-3.4,3.00) {serviceDiscovery\\[-2pt]{\tiny\rmfamily ServiceRegistry}};
\draw[thick,ArrowheadTeal,-{Stealth[length=1.6mm]}] (c0.east) -- (-1.8,3.00);
\node[ahcsm,anchor=east] (c1) at (-3.4,2.25) {identity\\[-2pt]{\tiny\rmfamily Authentication}};
\draw[thick,ArrowheadTeal,-{Stealth[length=1.6mm]}] (c1.east) -- (-1.8,2.25);
\node[ahcsm,anchor=east] (c2) at (-3.4,1.50) {authorization\\[-2pt]{\tiny\rmfamily ConsumerAuthorization}};
\draw[thick,ArrowheadTeal,-{Stealth[length=1.6mm]}] (c2.east) -- (-1.8,1.50);
\node[ahcsm,anchor=east] (c3) at (-3.4,0.75) {serviceOrchestration\\[-2pt]{\tiny\rmfamily DynamicServiceOrchestration}};
\draw[thick,ArrowheadTeal,-{Stealth[length=1.6mm]}] (c3.east) -- (-1.8,0.75);
\node[ahcsm,anchor=east] (c4) at (-3.4,0.00) {workflowWorkItem\\[-2pt]{\tiny\rmfamily WorkflowManager}};
\draw[thick,ArrowheadTeal,-{Stealth[length=1.6mm]}] (c4.east) -- (-1.8,0.00);
\node[ahcsm,anchor=east] (c5) at (-3.4,-0.75) {workflowContext\\[-2pt]{\tiny\rmfamily WorkflowManager}};
\draw[thick,ArrowheadTeal,-{Stealth[length=1.6mm]}] (c5.east) -- (-1.8,-0.75);
\node[ahcsm,anchor=east] (c6) at (-3.4,-1.50) {workflowLease\\[-2pt]{\tiny\rmfamily WorkflowManager}};
\draw[thick,ArrowheadTeal,-{Stealth[length=1.6mm]}] (c6.east) -- (-1.8,-1.50);
\node[ahcsm,anchor=east] (c7) at (-3.4,-2.25) {(step services)\\[-2pt]{\tiny\rmfamily Equipment / application systems}};
\draw[thick,ArrowheadTeal,-{Stealth[length=1.6mm]}] (c7.east) -- (-1.8,-2.25);
\node[ahcsm,anchor=east] (c8) at (-3.4,-3.00) {choreographerExecutor\\[-2pt]{\tiny\rmfamily Choreographer}};
\draw[thick,ArrowheadTeal,-{Stealth[length=1.6mm]}] (c8.east) -- (-1.8,-3.00);
\node[font=\scriptsize\bfseries,ArrowheadPurple] at (5.0,3.90) {produced};
\node[font=\scriptsize\bfseries,ArrowheadTeal] at (-5.0,3.90) {consumed};
\end{tikzpicture}
  }
  \caption{Block view of the WorkflowExecutor with its produced and consumed services.}
  \label{fig:block}
\end{figure}


\newpage
\section{Implementation}
\label{sec:implementation}

\subsection{Implementation language and tools}
Java~21 and Spring Boot~4 (Spring Framework~7, Spring Security~7), Spring Data JPA with Hibernate, Jackson for JSON,
Springdoc for the OpenAPI description, built with Maven~3.9 (in the \texttt{maven:3.\allowbreak{}9-\allowbreak{}eclipse-\allowbreak{}temurin-\allowbreak{}21} image). The
system builds on the vendored Arrowhead~5.4.0 common library (\texttt{vendor/\allowbreak{}common-\allowbreak{}utils},
\texttt{vendor/\allowbreak{}data-\allowbreak{}transfer-\allowbreak{}objects}): system registration, identity, TLS, management filters and HTTP consumption
of other services. The source is module \texttt{workflow-\allowbreak{}executor} of \texttt{Eclipse-\allowbreak{}Arrowhead-\allowbreak{}java/\allowbreak{}ah5-\allowbreak{}Workflow-\allowbreak{}java}; the
container image runs the Spring Boot jar on Eclipse Temurin~21.

\subsection{Internal structure}

{\small
\begin{longtable}{|p{6.2cm}|p{9.6cm}|}

\hline
\rowcolor{gray!33} Component & Responsibility \\ \hline
\endhead
{\footnotesize \texttt{Workflow\allowbreak{}Executor\allowbreak{}System\allowbreak{}Info}} & Registration of \texttt{workflow\allowbreak{}Execution} with metadata \texttt{workflow\allowbreak{}Capabilities}, \texttt{workflow\allowbreak{}Services}, \texttt{workflow\allowbreak{}Workstation}, \texttt{workflowMode} and \texttt{unrestricted\allowbreak{}Discovery}. \\ \hline
{\footnotesize \texttt{api.\allowbreak{}http.\allowbreak{}ExecutionAPI,\allowbreak{} GeneralManagementAPI}} & Service controllers; \texttt{start}/\texttt{abort} accept only the configured WorkflowManager or a sysop. \\ \hline
{\footnotesize \texttt{runtime.\allowbreak{}StepRunner}} & The execution of one work item (lease, orchestrate, compose, invoke, release, report) on a thread pool; abort and status. \\ \hline
{\footnotesize \texttt{runtime.\allowbreak{}PullLoop}} & Claims work items in pull mode. \\ \hline
{\footnotesize \texttt{runtime.\allowbreak{}ProviderResolver}} & Pull orchestration with the executor as requester at the configured orchestrator; static providers for development. \\ \hline
{\footnotesize \texttt{runtime.\allowbreak{}ManagerClient}} & Consumer of workflowWorkItem, workflowContext and workflowLease. \\ \hline
{\footnotesize \texttt{adapter.\allowbreak{}EquipmentAdapter,\allowbreak{} HttpJsonAdapter,\allowbreak{} Binding}} & Adapter SPI and the generic HTTP(S) JSON adapter. \\ \hline
{\footnotesize \texttt{api.\allowbreak{}compat.\allowbreak{}ChoreographerCompatAPI,\allowbreak{} ChoreographerCompatRunner,\allowbreak{} init.\allowbreak{}ChoreographerRegistrar}} & Choreographer executor interface; registration on start and unregistration on \texttt{Context\allowbreak{}Closed\allowbreak{}Event}. \\ \hline
{\footnotesize \texttt{security.\allowbreak{}ServiceAuthorizationFilter}} & Provider-side ConsumerAuthorization check. \\ \hline
\end{longtable}
}


\subsection{Functional properties implementation}
\texttt{StepRunner} implements the steps of the execution primitive in order: (1) for a lease
\texttt{reserve} step, reserve and poll until granted (the pending call keeps the step running); (2) for a lease
\texttt{use} step, renew with the token from $\xi(s_{res})$ and fail without a valid lease; (3) orchestrate
$\theta(s)$ with the executor as requester and report \texttt{BOUND} with the binding or \texttt{BIND\_FAILED}; the run
continues only when the WorkflowManager accepted the binding; (4) read $\rho$ and $\xi|dep(s)$, compose the input with
\texttt{InputComposer} and invoke through the first adapter that supports the binding's interface; (5) release a
lease if the step is a \texttt{release} step; store $\xi(s)$ with \texttt{put-result} and report \texttt{SUCCESS}, or
report \texttt{ERROR}/\texttt{FATAL\_ERROR}. A step whose service is \texttt{workflowLease} runs against the
WorkflowManager's lease service only. Service names are compared case-insensitively without separators, because the
Choreographer normalises definitions to camelCase.

\subsection{Non functional properties implementation}
\subsubsection{Security}
Spring Security filter chains of the common library: TLS client authentication, identity-token verification against
the Authentication system (outsourced authentication), the management filter and the workflow systems'
\texttt{Service\allowbreak{}Authorization\allowbreak{}Filter}, which maps method and path to service definition and operation and asks
ConsumerAuthorization (\texttt{authorization.\allowbreak{}verify}). Object-level checks (holder, owner, configured manager) are made in
the service logic and give 403.
\subsubsection{Power management}
No specific power management. Thread pools are bounded; the WorkflowManager's watchdog and the executor's pull loop are
the only periodic activities.
\subsubsection{Internal monitoring}
Spring Boot health endpoint (\texttt{Workflow\allowbreak{}Executor\allowbreak{}Health\allowbreak{}Indicator}),
the common-library log table exposed through \texttt{general\allowbreak{}Management.\allowbreak{}get-\allowbreak{}log}, and optional sysMonitor telemetry over
AMQP (\texttt{sysmonitor.\allowbreak{}enabled}).
\subsubsection{Configuration}

{\small
\begin{longtable}{|p{5.6cm}|p{3.4cm}|p{6.8cm}|}

\hline
\rowcolor{gray!33} Property & Default & Meaning \\ \hline
\endhead
\texttt{server.\allowbreak{}port} & 8483 & Arrowhead connector port. \\ \hline
\texttt{workflow.\allowbreak{}executor.\allowbreak{}capabilities} & (empty) & $cap(w)$, comma separated. \\ \hline
\texttt{workflow.\allowbreak{}executor.\allowbreak{}services} & * & Services the executor can run. \\ \hline
\texttt{workflow.\allowbreak{}executor.\allowbreak{}workstation} & (empty) & Workstation. \\ \hline
\texttt{workflow.\allowbreak{}executor.\allowbreak{}mode} & push & \texttt{push} or \texttt{pull}. \\ \hline
\texttt{workflow.\allowbreak{}executor.\allowbreak{}orchestrator} & Dynamic\allowbreak{}Service\allowbreak{}Orchestration & The orchestrator asked for providers. \\ \hline
\texttt{workflow.\allowbreak{}executor.\allowbreak{}unrestricted-\allowbreak{}discovery} & true & Register \texttt{workflow\allowbreak{}Execution} with \texttt{unrestricted\allowbreak{}Discovery}, so the WorkflowManager finds it under RESTRICTED discovery. \\ \hline
\texttt{workflow.\allowbreak{}executor.\allowbreak{}poll-\allowbreak{}ms} & 500 & Pull interval. \\ \hline
\texttt{workflow.\allowbreak{}executor.\allowbreak{}threads} & 8 & Work item threads. \\ \hline
\texttt{workflow.\allowbreak{}executor.\allowbreak{}lease-\allowbreak{}poll-\allowbreak{}ms} & 200 & Lease poll interval while queued. \\ \hline
\texttt{workflow.\allowbreak{}manager.\allowbreak{}system-\allowbreak{}name} & WorkflowManager & The only system allowed to start/abort work items. \\ \hline
\texttt{workflow.\allowbreak{}manager.\allowbreak{}url} & (empty) & Direct WorkflowManager URL for development; otherwise discovered. \\ \hline
\texttt{workflow.\allowbreak{}executor.\allowbreak{}static-\allowbreak{}providers} & (empty) & Static providers for development (no orchestration). \\ \hline
\texttt{workflow.\allowbreak{}choreographer.\allowbreak{}enabled} & false & Choreographer executor compatibility. \\ \hline
\texttt{workflow.\allowbreak{}choreographer.\allowbreak{}system-\allowbreak{}name} & Choreographer & Choreographer allowed to call the compatibility interface. \\ \hline
\texttt{workflow.\allowbreak{}choreographer.\allowbreak{}min-\allowbreak{}version /\allowbreak{} max-\allowbreak{}version} & 1 / 1 & Versions reported in Choreographer \texttt{service-info}. \\ \hline
\end{longtable}
}


\newpage
\section{Services}
\label{sec:services}
\subsection{Produced services}
All produced services use interface template \texttt{generic\_\allowbreak{}https}, policy \texttt{CERT\_AUTH} and version 1.0.0 and are
registered with every operation's method and path.

{\small
\begin{longtable}{|p{3.8cm}|p{6.2cm}|p{5.8cm}|}

\hline
\rowcolor{gray!33} Service & Operations & Documents \\ \hline
\endhead
\texttt{workflow\allowbreak{}Execution} & service-info, start, abort, get-status & \texttt{Workflow\allowbreak{}Executor-\allowbreak{}v5.\allowbreak{}4.\allowbreak{}0-\allowbreak{}workflow\allowbreak{}Execution-\allowbreak{}SD}, \texttt{-HTTP-IDD} \\ \hline
\texttt{general\allowbreak{}Management} & get-log, get-config & \texttt{Workflow-\allowbreak{}v5.\allowbreak{}4.\allowbreak{}0-\allowbreak{}general\allowbreak{}Management-\allowbreak{}SD}, \texttt{-HTTP-IDD} \\ \hline
\texttt{Choreographer executor interface} & service-info, start, abort (not registered; under \textbackslash{}texttt\{/workflowexecutor/choreographer\}) & Choreographer executor IDD (Choreographer 5.x) \\ \hline
\end{longtable}
}

\subsection{Consumed services}

{\small
\begin{longtable}{|p{3.4cm}|p{3.2cm}|p{3.6cm}|p{5.6cm}|}

\hline
\rowcolor{gray!33} Service & Provider & Operations & Purpose \\ \hline
\endhead
\texttt{service\allowbreak{}Discovery} & ServiceRegistry & register, revoke & Registration of \texttt{workflow\allowbreak{}Execution} with capability metadata. \\ \hline
\texttt{identity} & Authentication & login, logout, identity-verify & Own identity token; verification of callers. \\ \hline
\texttt{authorization} & ConsumerAuthorization & verify & Provider-side authorization. \\ \hline
\texttt{service\allowbreak{}Orchestration} & Dynamic\allowbreak{}Service\allowbreak{}Orchestration & pull & Provider of each step, with the executor as requester. \\ \hline
\texttt{workflow\allowbreak{}Work\allowbreak{}Item} & WorkflowManager & claim, report & Pull and reports. \\ \hline
\texttt{workflow\allowbreak{}Context} & WorkflowManager & get-parameters, get-results, put-result & $\rho$, $\xi|dep(s)$, $\xi(s)$. \\ \hline
\texttt{workflowLease} & WorkflowManager & reserve, poll, renew, release & Shared resources. \\ \hline
(step services) & Equipment / application systems & as orchestrated & The work itself. \\ \hline
\texttt{choreographer\allowbreak{}Executor} & Choreographer & register, unregister, notify-step-done & Only with \texttt{workflow.\allowbreak{}choreographer.\allowbreak{}enabled}. \\ \hline
\end{longtable}
}


\section{Verification}
\label{sec:verification}
Seven unit tests with fake equipment (passed on 2026-11-24), and the smoke and cloud tests of the WorkflowManager, in which three executors
run push and pull sessions and the plant sees the executors' own certificates. Choreographer compatibility was
verified against the cloud's Choreographer 5.4.0: a compiled recipe ran as a Choreographer plan on an executor that
registered itself, reported with \texttt{notify-\allowbreak{}step-\allowbreak{}done} and unregistered on shutdown (34/34 checks).


\section{Implementation variants}
\label{sec:variants}
\begin{itemize}
\item \textbf{Java (reference implementation).} \texttt{Eclipse-\allowbreak{}Arrowhead-\allowbreak{}java/\allowbreak{}ah5-\allowbreak{}Workflow-\allowbreak{}java/\allowbreak{}workflow-\allowbreak{}executor}, Java 21, Spring
Boot, built with Maven together with the shared \texttt{workflow-\allowbreak{}model} module and the vendored Arrowhead~5.4.0 common
library.
\item \textbf{Go.} Not implemented. No workflow system exists in \texttt{Eclipse-\allowbreak{}Arrowhead-\allowbreak{}go}.
\item \textbf{Rust.} Not implemented. No workflow system exists in \texttt{Eclipse-\allowbreak{}Arrowhead-\allowbreak{}rust}.
\end{itemize}
The workflow systems are therefore not yet interoperable across the three implementations. A Go or Rust port must
reproduce the IDDs of this set exactly and, for the WorkflowManager, the behaviour of \texttt{workflow-\allowbreak{}model}: the
recipe schema and decoder, the validator's findings and verdicts, the guard grammar and the transition rules. The
architecture prescribes conformance vectors (recipe in, verdict and transition trace out) to be run by all three
implementations; they are not written yet.


\newpage
\bibliographystyle{IEEEtran}
\bibliography{bibliography}

\newpage

\section{Revision History}
\subsection{Amendments}

\noindent\begin{tabularx}{\textwidth}{| p{1cm} | p{2.2cm} | p{1.6cm} | X | p{3.6cm} |} \hline
\rowcolor{gray!33} No. & Date & Version & Subject of Amendments & Author \\ \hline
1 & 2026-11-23 & \arrowversion & Initial edition: system design description of the WorkflowExecutor, written from the Java reference implementation (ah5-Workflow-java). & Documentation maintenance \\ \hline
2 & 2026-11-24 & \arrowversion & Verification status. & Documentation maintenance \\ \hline
\end{tabularx}

\subsection{Quality Assurance}

\noindent\begin{tabularx}{\textwidth}{| p{1cm} | p{2.2cm} | p{1.6cm} | X |} \hline
\rowcolor{gray!33} No. & Date & Version & Approved by \\ \hline
1 & & \arrowversion & Pending review \\ \hline
\end{tabularx}

\end{document}
